\section{Analysis Protocol and Experimental Setup}
\label{sec:method}

\texttodo{Add in the correlation with our oracle, the distance function and the advantages for all algorithms}

We study how offline GCRL value-learning objectives interact with behavioral policy extraction through phase-resolved \emph{trainability landscapes}.
The matched hyperparameter space is \(\Lambda=\Lambda_{\eta}\times\Lambda_{\alpha}\), where \(\eta\) is the optimizer learning rate and \(\alpha\) is the AWR temperature.
Thus, \(\eta\) primarily affects how the value, distance, or reachability signal is learned, while \(\alpha\) controls how strongly the actor-facing advantage is extracted into a policy.

\paragraph{Datasets and phased response surfaces.}
Our main experiments use static OGBench datasets, focusing on AntMaze for long-horizon locomotion and Cube/Scene for manipulation-oriented goal reaching.
Following \citet{mohan-automlconf23a}, training is divided into phases.
At each phase boundary \(t\), all configurations \(\lambda\in C\subset\Lambda_{\eta}\times\Lambda_{\alpha}\) are evaluated, yielding a response surface \(f_t:C\to\mathbb{R}\).
The next phase is initialized from a shared reference checkpoint selected from the previous phase, so within-phase comparisons are not confounded by incomparable training histories.
We visualize response surfaces with Independent Gaussian Process Regressors (IGPRs), but compute all quantitative metrics from raw evaluations.

\paragraph{Evaluation and landscape breadth.}
We evaluate downstream behavior with OGBench's success rate and normalized goal-distance return,
\[
    R_{\mathrm{dist}}(s_T,g)=1-\|s_T-g\|_2/\|s_0-g\|_2.
\]
Success rate is used for checkpoint selection, while \(R_{\mathrm{dist}}\) is used for landscape visualization and geometry analysis.
We do not treat \(R_{\mathrm{dist}}\) as an oracle for advantage quality, since the algorithms expose actor-facing advantages with different semantics.
To measure landscape breadth, we compute the \(\varepsilon\)-optimality mass
\[
    \rho_{\varepsilon}(t) = |C|^{-1} \left| \left\{ \lambda\in C: f_t(\lambda)\ge \varepsilon\max_{\lambda'\in C}f_t(\lambda') \right\} \right|.
\]
We report \(\rho_{0.9}\) and \(\rho_{0.8}\), which respectively measure robustness of the near-peak region and the broader moderate-performance basin.

\paragraph{AWR concentration diagnostics.}
To explain why some landscapes are broad while others are sharp, we inspect the actor-side distribution induced by AWR.
For a logged batch, let \(A_i=A_{\mathrm{alg}}(s_i,a_i,g_i)\), \(z_i=\alpha A_i\), and \(w_i=\min\{\exp(z_i),w_{\max}\}\).
We partition samples into downweighted, active-unsaturated, and saturated regions using \(z_i<0\), \(0\le z_i\le \log w_{\max}\), and \(z_i>\log w_{\max}\), and track both sample fractions \((p_-,p_{\mathrm{mid}},p_{\mathrm{sat}})\) and normalized weight masses \((W_-,W_{\mathrm{mid}},W_{\mathrm{sat}})\).
We also measure the normalized effective sample size,
\[
    \mathrm{ESS} = \frac{(\sum_i w_i)^2}{B\sum_i w_i^2},
\]
where high ESS indicates diffuse weighted cloning and low ESS indicates concentrated actor updates.
We interpret ESS jointly with return and advantage semantics, since concentration is useful only when the selected samples are reliable actor-facing winners.

\paragraph{Future-vs-random advantage diagnostics.}
To compare advantage semantics across methods, we use a fixed diagnostic goal table.
For each anchor transition \((s_i,a_i)\), we sample one future goal \(g_i^+\sim p_{\mathcal D}^{\mathrm{geom}}(\cdot\mid s_i,a_i)\) and \(K=32\) random goals \(g_{ij}^-\sim p_{\mathcal D}^{\mathrm{rand}}\), reusing the same table across algorithms, configurations, seeds, and phases whenever possible.
We compute
\[
    \Delta A_{ij}^{\mathrm{FR}} = A_{\mathrm{alg}}(s_i,a_i,g_i^+) - A_{\mathrm{alg}}(s_i,a_i,g_{ij}^-).
\]
This tests whether the actor-facing advantage ranks a dataset-reachable future goal above a random goal for the same anchor transition.
Before clipping, AWR amplifies this gap through
\[
    w_{\alpha}(s_i,a_i,g_i^+)/ w_{\alpha}(s_i,a_i,g_{ij}^-) = \exp(\alpha\Delta A_{ij}^{\mathrm{FR}}).
\]
We summarize future-vs-random separation with \(\mathrm{FR\text{-}AUC}=\Pr(\Delta A^{\mathrm{FR}}>0)\) and
\[
    \mathrm{EI} = \mathbb{E}[\Delta A^{\mathrm{FR}}]/ \bigl(\mathrm{Std}(\Delta A^{\mathrm{FR}})+\epsilon\bigr).
\]
FR-AUC measures whether the gap has the correct sign more often than chance, while EI measures the signal-to-noise ratio of the gap.
